\chapter{Introdução}
\label{cha:introducao}

O metaverso é um conceito multidimensional cuja definição permanece objeto de debate na literatura científica. Revisões recentes tratam o tema por perspectivas diferentes: algumas concentram-se na definição conceitual; outras investigam tecnologias, aplicações, arquiteturas ou interoperabilidade. Essa diversidade torna inadequado tratar uma única formulação como definição universal.

Ritterbusch e Teichmann (2023) realizaram uma revisão sistemática especificamente dedicada à definição do metaverso e analisaram 28 definições e descrições científicas. O estudo mostra a variedade de formulações utilizadas na literatura e constitui uma referência para a análise conceitual realizada neste trabalho. \cite{ritterbusch2023}

A literatura posterior amplia o problema. Tukur et al. (2024) realizaram uma revisão de escopo sobre técnicas, tecnologias e aplicações do metaverso, identificando 12 classes tecnológicas e aplicações em 11 setores. Boztas, Ghadafi e Ibrahim (2025), por sua vez, realizaram uma revisão sistemática de 103 estudos, dos quais 33 apresentavam definições explícitas e 25 propunham modelos arquiteturais. \cite{tukur2024,boztas2025}

Diante desse cenário, o presente trabalho não pretende declarar uma nova definição universal. Seu objetivo é comparar a definição adotada no projeto com perspectivas científicas verificáveis e, a partir dessa comparação, construir uma definição operacional adequada ao escopo do projeto.

A questão central é, portanto, distinguir três níveis que frequentemente aparecem misturados: as propriedades que caracterizam o fenômeno, as tecnologias que podem sustentá-lo e as aplicações que podem ser desenvolvidas sobre ele. Essa separação é necessária para evitar que tecnologias específicas, como realidade virtual, inteligência artificial ou blockchain, sejam tratadas automaticamente como requisitos universais.

O objetivo geral é analisar comparativamente a literatura científica sobre o metaverso e confrontá-la com a definição de trabalho adotada no projeto. Como objetivos específicos, busca-se: identificar características recorrentes nas definições; comparar definições e arquiteturas recentes; distinguir propriedades conceituais de tecnologias habilitadoras; analisar aplicações e interoperabilidade; e formular uma definição operacional para o projeto.

A contribuição deste trabalho é, portanto, uma síntese comparativa fundamentada em fontes científicas identificáveis. Sempre que um número é apresentado, ele é atribuído à fonte que o produziu. Não são apresentados como resultados desta pesquisa números que não tenham sido efetivamente coletados e registrados pelo presente estudo.

O restante do trabalho está organizado da seguinte forma: o Capítulo~\ref{cha:revisao} apresenta a revisão da literatura; o Capítulo~\ref{cha:questoes} apresenta as questões de pesquisa; o Capítulo~\ref{cha:relacionados} discute os trabalhos relacionados; o Capítulo~\ref{cha:metodo} descreve o método e a rastreabilidade das evidências; o Capítulo~\ref{cha:critica} apresenta a análise crítica e a definição operacional proposta; o Capítulo~\ref{cha:resultados} sintetiza os resultados; e o Capítulo~\ref{cha:conclusao} apresenta as conclusões e limitações.
